\section{The objective: more proof construction, not a bigger timeout}

\forge{} is a \emph{proof-producing automation system with a structural search
layer}. It is not a new foundational logic, not a replacement for Lean's
kernel, and not a rewrite of \grind{}. It preserves the existing kernel and
treats search engines, learned rankings, computer-algebra systems, and
numerical solvers as fallible advisers whose output must be checked before it
becomes a fact.

The design separates two tasks that are commonly conflated. A
\emph{structural planner} chooses a proof shape: introduce a binder,
construct a witness, apply a theorem, prove an auxiliary lemma, generalise an
accumulator, or use an induction principle. A \emph{local saturation engine}
derives consequences within the resulting context. \grind{} is an excellent
candidate for the second role. The proposed improvement is to supply proof
structure and certified facts that local saturation cannot obtain on its own,
and to feed useful consequences back rather than restarting every tool from an
unchanged goal.

\[
\boxed{\text{local consequence search}}
\quad\longleftrightarrow\quad
\boxed{\text{structural proof search}}
\]

Three short examples fix the pattern. A nonlinear worker may prove
$0\le p(x,y)$, combine it with an existing $p(x,y)\le0$, and obtain
$p(x,y)=0$; that equality is then useful to congruence closure even when the
final target is an equality between uninterpreted function applications. An
induction planner may generalise an accumulator, synthesise its polynomial
invariant, and invoke \grind{} only on the base and step obligations. A
witness worker may construct a parametric solution of a linear system and
leave integrality or bound obligations to other workers. In all three cases
the benefit comes from \emph{coordination around checked intermediate claims},
not from any single engine becoming stronger.

\subsection{Two meanings of ``more powerful''}

The desired outcome is substantially higher automatic coverage on useful Lean
developments: \emph{more goals solved automatically from comparable initial
information}, particularly goals that need structural proof steps or new
terms. That is not the same as a stronger logic, a change to Lean's trusted
axioms, uniform speed improvement, logical completeness, or success-set
inclusion at a fixed total budget. Those are different claims and they should
never be argued for interchangeably.

\begin{definition}
For a tactic implementation $T$, an environment $E$, a goal $G$, and a
resource allocation $B$, write $T(E,G;B)\downarrow$ when $T$ returns a proof
accepted under the selected trust policy within that allocation.
\end{definition}

\forge{} offers two modes, and an evaluation must report both.

In \emph{compatibility mode} it first runs the stock \grind{} invocation, with
the same input, options, and solver budget, before running any new machinery.
In \emph{fixed-budget mode} it partitions a single budget among workers and
may sacrifice some baseline successes in exchange for others.

\begin{proposition}[Conditional baseline preservation]
\label{prop:preservation}
Assume the stock run is reproduced deterministically with identical inputs and
resources, and that its successful proof is returned immediately. If
compatibility mode allocates that full run before spending any additional
extension resources, then every success of the stock invocation is also a
success of compatibility mode.
\end{proposition}
\begin{proof}
The first stage is the stock invocation itself. A success exits before any
extension can change its state or consume its reserved allocation.
\end{proof}

Proposition~\ref{prop:preservation} is a \emph{scheduling} property, not a
superiority result at equal cost. Process startup, memory pressure, imports,
and global wall-clock limits can invalidate an operational reading of it
unless they are controlled. Nor can any finite-budget heuristic promise
pointwise dominance over another heuristic at exactly the same cost. A fair
experiment reports compatibility mode and fixed-budget mode separately, and
never presents the preservation property as an equal-cost performance theorem.

Failure-triggered intervention has recent precedent in research on \grind{}
itself; the cascade is prior art and is not claimed here as an
invention~\cite{learned-interventions}.

\subsection{Where new coverage can realistically come from}

The target is not ``everything \grind{} cannot prove''. It is a specific
collection of mechanisms with a plausible cost/coverage tradeoff:

\begin{enumerate}[leftmargin=1.4em]
\item \textbf{Proof structure.} Select induction principles, generalise
  changing state, construct existential witnesses, and establish auxiliary
  lemmas before using them.
\item \textbf{Missing algebraic consequences.} Discover hidden sums of
  squares, domain-dependent polynomial positivity, and exact zero information
  that can be shared with other engines.
\item \textbf{Directed obligation solving.} Choose theorem applications from
  the conclusion backward, create their side conditions explicitly, and use
  successful side conditions to activate guarded forward reasoning.
\end{enumerate}

These are proposed extensions, not assertions that every example requiring one
will fail under all existing \grind{} configurations or library annotations.
Some will already be solved by \grind{}, \code{nlinarith}, Aesop, or a
suitable library theorem. Only controlled runs on a real corpus can identify
genuinely new coverage.

\subsection{The deliverable and its evidence levels}

This document deliberately keeps four things apart, and the reader is asked to
keep them apart too:

\begin{enumerate}[leftmargin=1.4em]
\item \textbf{Mathematical soundness arguments.} The propositions and theorems
  below explain why each certificate format implies the statement it claims.
  These are ordinary mathematics and are independent of any implementation.
\item \textbf{Tested Python checking.} Executable prototypes synthesise and
  re-check certificates in exact rational arithmetic, with negative and
  mutation tests. This is evidence about programs, not about Lean.
\item \textbf{Generated Lean candidate source.} Some workers emit Lean proof
  scripts. Where a script has not been compiled, it is \emph{candidate
  source}, and saying so is not a formality.
\item \textbf{The proposed production implementation.} The Lean tactic, its
  reflected checkers, its \grind{} adapter, and its measured performance.
  None of this exists yet.
\end{enumerate}

\begin{remark}
A Python check is not a Lean-kernel proof. The prototype checkers are
hand-written research code; several share representation and matching routines
with the searches they are meant to audit. ``Search-free replay'' means the
numerical and symbolic search libraries are not loaded --- it does not mean an
independently implemented formal verification.
\end{remark}

\subsection{What was actually built, and what was not}

Across the nine contributing efforts, the following algorithm families were
implemented and executed in Python: exact sparse rational polynomial
arithmetic; exact quadratic decomposition by symmetric pivoting; finite
polynomial-cone search with rational repair; Bernstein box subdivision with
checked coverage; square-factor and Sturm certificates for univariate
intervals; polynomial recurrence and accumulator-invariant synthesis by exact
CEGIS; affine and integer-lattice witness synthesis; demand-directed ground
Horn reasoning; a typed equational prover with structural-induction replay;
proof-logging CDCL with an integer-difference-logic theory solver; and
standard-library-only replay of every stored certificate.

The following were \emph{not} built by anyone: an installed \code{forge}
tactic; a Lean formalisation of any checker; the full obligation scheduler
against real Lean snapshots; native \grind{} state access; a dependent-type
reifier; general quantifier instantiation; higher-order synthesis; external
SMT proof reconstruction; a semidefinite backend; and any measured comparison
against a Lean tactic.

Section~\ref{sec:experiments} reports what each of the eighteen runs actually
observed. Those runs used different seeds, different generators, and different
case sets; their counts are not comparable and must never be pooled into a
single accuracy figure.
